\documentclass{article}
\usepackage{geometry}
\geometry{a4paper, margin=1in}
\usepackage{graphicx}
\usepackage{amsmath}
\usepackage{amsfonts}
\usepackage{amssymb}
\usepackage{eufrak}
\usepackage{physics}
\usepackage{hyperref}

\title{Harmonic Hyperbolic Networks (H$^2$N)}
\author{Alfred Ricker}
\date{December 2025}

\newcommand{\D}{\mathrm{d}}
\newcommand{\mathbbD}{\mathbb{D}}
\newcommand{\boundary}{\partial \mathbb{D}}

\begin{document}

\maketitle

\subsection{Warning}
This outline was almost entirely generated by Gemini as a sketch, but it doesn't prove very useful to me because I have a hard time understanding it. It also (very likely) suffers from painfully unstable gradients and raises many information synthesis problems.

\section{Introduction}
The Harmonic Hyperbolic Network (H$^2$N) represents a fundamental departure from discrete, digital neural processing towards a continuous, resonant dynamical system. Rather than modeling intelligence as the sequential accumulation of state bits, H$^2$N models it as the interference of complex waves within a non-Euclidean geometry. This framework unifies the hierarchical nature of knowledge with the spectral nature of signal processing.

\section{Desirable Properties of a Learning Algorithm}
Before deriving the formalism, we establish the core philosophical and biological constraints that guide the design of the H$^2$N architecture.

\subsection{Hierarchy and Zipf's Law}
Human knowledge is not flat; it is deeply hierarchical. General concepts (e.g., ``Mammal'') are few but highly connected, while specific instances (e.g., ``The stray cat I saw yesterday'') are numerous but sparsely connected. A learning system should naturally embed this structure. Euclidean space fails here, as volume grows polynomially ($r^n$). We require a space where volume grows exponentially ($e^r$) to accommodate the exponential explosion of specificities relative to generalities, matching the statistical distribution of language (Zipf's Law).

\subsection{Continuum and Resonance}
Biological inputs (sound, light, motor control) are continuous signals, yet traditional ANNs discretize them into static vectors. A more natural approach is to treat information as \textit{waves}. This allows for:
\begin{itemize}
    \item \textbf{Superposition:} Multiple concepts can coexist in the same space without corruption.
    \item \textbf{Interference:} Contradictory ideas can physically cancel each other out.
    \item \textbf{Resonance:} The system can "tune" itself to specific frequencies, allowing memory to be a passive geometric property rather than an active storage cost.
\end{itemize}

\subsection{Locality and Plasticity}
In a biological brain, there is no central processor calculating a global loss gradient. Updates must be \textit{local}. A neuron should learn solely based on the signals it receives from its immediate neighbors and the feedback "echo" from the environment. Furthermore, the system must balance stability (memory) and plasticity (learning). This requires a physical mechanism where neurons can "freeze" successful configurations and "melt" erroneous ones.

\section{Polarization and Hyperbolicity}
To satisfy the requirements of hierarchy and complex interference, we separate the network's geometry into two distinct manifolds: the \textit{Base Space} and the \textit{Fiber Space}.

\subsection{The Base: Hyperbolic Space ($\mathbb{D}$)}
The physical location of a neuron lives in the Poincaré disk $\mathbb{D}$ (typically 2D or 3D). This low dimensionality is sufficient to define the \textit{hierarchy}.
\begin{itemize}
    \item \textbf{Radial Depth ($|z|$):} Encodes abstraction. $|z| \approx 0$ represents fundamental concepts; $|z| \to 1$ represents sensory details.
    \item \textbf{Distance ($d_{\mathbb{D}}$):} Encodes latency and decay. It acts as the medium that attenuates signals, ensuring that only logically related concepts (those close in space) interact strongly.
\end{itemize}

\subsection{The Fiber: Spinor Space ($\mathbb{C}^N$)}
If we only used scalar waves, "Apple" and "Orange" would interfere (add or subtract). To allow orthogonal concepts to pass through each other, the signal itself must live in a higher-dimensional \textit{Polarization Space}.
Each point $z \in \mathbb{D}$ is equipped with a vector space $\mathbb{C}^N$ (a spinor bundle). Information is encoded in the \textit{orientation} of the wave in this high-dimensional fiber. This allows the network to process $N$ independent features simultaneously within the same physical region.

\section{State Variables}
Each computational unit (neuron) $j$ is described by the following set of dynamical variables:

\begin{description}
    \item[Position $z_j \in \mathbb{D}$:] The coordinate in the hyperbolic hierarchy. Determines frequency $\omega(|z|)$ and connectivity.
    \item[Gain $A_j(t) \in \mathbb{R}^+$:] The metabolic energy of the node. High gain implies the node is currently relevant and acting as a repeater; low gain implies quiescence.
    \item[Spinor State $\mathbf{s}_j(t) \in \mathbb{C}^N$:] The instantaneous complex state vector, encoding phase and polarization.
    \item[Plasticity $\eta_j(t)$ :] A coefficient of freedom between 0 and 1. Determines how easily the node can move or change its internal polarization.
\end{description}

\section{Wave State and Propagation}

\subsection{The Spinor Bundle and Polarization}
In scalar field theories, wave interference is binary: signals either constructively reinforce or destructively cancel. This is insufficient for a cognitive architecture, where orthogonal concepts (e.g., ``Color'' and ``Velocity'') must coexist in the same spatial region without corruption.

To resolve this, we define the state of the network not as a scalar field, but as a section of a \textbf{Spinor Bundle} $S$ over the manifold $\mathbb{D}$. At every point $z$, the wave function $\mathbf{\Psi}(z, t)$ lives in an internal vector space $\mathbb{C}^N$:
\begin{equation}
    \mathbf{\Psi}(z, t) \in \Gamma(S), \quad \mathbf{\Psi}(z, t) = \sum_{\mu=1}^N \psi^\mu(z, t) \mathbf{e}_\mu
\end{equation}
where $\{\mathbf{e}_\mu\}$ forms the orthonormal basis of the fiber.



This introduces the degree of freedom of \textit{Polarization}. The interaction between two intersecting waves is governed by the Hermitian inner product of their spinors:
\begin{equation}
    \mathcal{I}_{int} = \text{Re} \langle \mathbf{\Psi}_1, \mathbf{\Psi}_2 \rangle = |\mathbf{\Psi}_1| |\mathbf{\Psi}_2| \cos(\Delta \phi) \cos(\theta_{spin})
\end{equation}
Here, $\theta_{spin}$ represents the angle in the high-dimensional feature space. This allows for \textit{Selective Interference}:
\begin{itemize}
    \item \textbf{Parallel ($\theta=0$):} Full interference (synonyms/antonyms).
    \item \textbf{Orthogonal ($\theta=\pi/2$):} Zero interference. Distinct concepts pass through each other transparently, allowing the network to multiplex different streams of thought.
\end{itemize}

\subsection{The Hyperbolic Propagator (Green's Function)}
A neuron $j$ at position $z_j$ acts as a point source radiator. The propagation of its signal through the bulk is governed by the Green's function $G(z, z_j)$ of the damped wave equation on $\mathbb{H}^n$.

Unlike Euclidean space where flux creates a $1/r^{n-2}$ decay, the exponential volume expansion of hyperbolic space causes rapid geometric attenuation. We approximate the propagator in the high-frequency limit as:
\begin{equation}
    G(z, z_j; t) \approx \underbrace{\frac{1}{\left(\sinh d_{\mathbb{D}}(z, z_j)\right)^{(D-1)/2}}}_{\text{Geometric Decay}} \cdot \underbrace{e^{i(k d_{\mathbb{D}}(z, z_j) - \omega_j t)}}_{\text{Phase Evolution}} \cdot \underbrace{e^{-\mu d_{\mathbb{D}}(z, z_j)}}_{\text{Mass Damping}}
\end{equation}

\begin{itemize}
    \item The $\sinh$ term dominates the denominator, ensuring locality. A neuron's influence is effectively bounded to its local neighborhood unless highly amplified.
    \item The Mass Damping $\mu$ is a tunable hyperparameter that defines the ``horizon'' of a concept's influence.
\end{itemize}

\subsection{Discrete-to-Continuum Coupling: Ricker Wavelets}
Inputs to the system are discrete events (tokens). To map these into the continuous dynamical system without introducing infinite-frequency artifacts (which occur with Dirac deltas), we utilize \textit{Ricker Wavelets} (the second derivative of a Gaussian).

The input field $\mathbf{I}_{in}(z, t)$ is constructed by convolving the spatiotemporal token locations with the Ricker kernel:
\begin{equation}
    \mathbf{I}_{in}(z, t) = \sum_{k} \delta(z - z_k) \cdot \mathbf{\hat{u}}_k \cdot \left( 1 - \omega_0^2 (t - t_k)^2 \right) e^{-\frac{1}{2}\omega_0^2 (t - t_k)^2}
\end{equation}

This wavelet has two critical properties for neural encoding:
\begin{enumerate}
    \item \textbf{Zero Mean:} $\int_{-\infty}^{\infty} W(t) dt = 0$. This ensures inputs do not introduce a DC bias that would destabilize the total energy of the system.
    \item \textbf{Spectral Purity:} It is well-localized in both time and frequency, minimizing spectral leakage between adjacent concepts in the hierarchy.
\end{enumerate}

\subsection{Holographic Renormalization and Boundary Integrals}
A central challenge in hyperbolic geometry is that any signal originating in the bulk decays exponentially as $e^{-(D-1)r/2}$ towards the boundary $\boundary$. A naive sensor at the boundary would measure zero signal.

To recover the data, we appeal to the AdS/CFT correspondence principle of \textit{Holographic Renormalization}. We assert that the physical observable is not the raw field $\mathbf{\Psi}$, but the \textit{source term} of the field's asymptotic expansion.

Near the boundary ($r \to 1$), the field behaves as:
\begin{equation}
    \mathbf{\Psi}(r, \theta) \sim A(\theta)(1-r)^{\Delta_-} + B(\theta)(1-r)^{\Delta_+}
\end{equation}
We isolate the non-decaying component by applying a counter-term scaling factor $\Omega(r)^{-\Delta}$, where $\Omega$ is the conformal factor of the metric. The renormalized boundary integral is:
\begin{equation}
    \mathbf{O}(\theta, t) = \lim_{r \to 1} \left( \frac{2}{1-r^2} \right)^{\Delta} \int_{\mathcal{N}} \mathbf{\Psi}_j(r, \theta, t) \, \D \mu(z_j)
\end{equation}
This operation effectively ``peels off'' the geometric curvature, projecting the bulk interference pattern onto the flat boundary screen. This ensures that abstract concepts deep in the core ($z \approx 0$) have the same causal authority at the output as sensory details near the edge.

\section{Equations of State}
The system evolves according to a set of coupled differential equations governing movement, gain, and learning.

\subsection{Gain Dynamics (Metabolism)}
Nodes amplify signals that resonate with the solution. We define an ``Error Echo'' field $\mathbf{\Psi}^\dagger(z)$ which propagates the boundary error backwards into the bulk.
\begin{equation}
    \tau_A \frac{\D A_j}{\D t} = \underbrace{-A_j}_{\text{Decay}} + \underbrace{\mu \tanh\left( \text{Re}\langle \mathbf{\Psi}_j, \mathbf{\Psi}^\dagger_j \rangle \right)}_{\text{Utility Gain}}
\end{equation}
If the node's forward state $\mathbf{\Psi}_j$ correlates with the backward echo $\mathbf{\Psi}^\dagger_j$, the node is constructive, and its gain increases.

\subsection{Plasticity (Annealing)}
Plasticity is modulated by the global reward signal $R(t)$. Successful outcomes freeze the network; failures melt it.
\begin{equation}
    \frac{\D \eta_j}{\D t} = -\gamma R(t) \cdot \eta_j \quad \text{(subject to } \eta \in [\eta_{min}, 1]\text{)}
\end{equation}

\subsection{Positional Dynamics (The Tuning Rule)}
The primary learning mechanism is the physical migration of nodes to resonant locations. The velocity of a node is determined by:
\begin{equation}
    \frac{\D z_j}{\D t} = \eta_j \left[ \underbrace{-\alpha z_j}_{\text{Zipf Force}} + \underbrace{\sum_k \frac{z_j - z_k}{|z_j - z_k|^2}}_{\text{Structural Repulsion}} + \underbrace{\beta R(t) \nabla_{z_j} \text{Re}\langle \mathbf{\Psi}_j, \mathbf{\Psi}^\dagger_j \rangle}_{\text{Adjoint Gradient}} \right]
\end{equation}
The final term is the \textit{Holographic Backpropagation} term. It moves the node along the gradient of the interference pattern, effectively tuning the phase delay to maximize constructive interference with the target output.

\end{document}